% =============================================================================
% capitolo3_implementazione.tex
% Capitolo 3: Architettura e Dettagli Implementativi
%
% LINEE GUIDA VITALI:
%  Questo e' il capitolo tecnico per eccellenza: il "COME".
%  Il lettore deve poter replicare il sistema dopo averlo letto.
%
%  A) ARCHITETTURA A BASSO LIVELLO:
%     Diagramma delle componenti con dipendenze.
%     Struttura delle directory/moduli principali.
%     Scelte architetturali (microservizi vs monolite, REST vs GraphQL, ecc.)
%     giustificate rispetto alle alternative scartate.
%
%  B) TECNOLOGIE E DIPENDENZE:
%     Linguaggi, framework, librerie. Per ognuna: versione, ruolo, motivo della scelta.
%
%  C) FILE / MODULI IMPORTANTI:
%     Descrivi i file/classi/moduli piu' rilevanti.
%     Usa \begin{lstlisting}...\end{lstlisting} per i code snippet.
%
%  D) ALGORITMI CHIAVE:
%     Spiega in dettaglio gli algoritmi non banali che hai implementato.
%     Pseudocodice, flowchart, o codice reale con spiegazione passo-passo.
%
%  E) MODELLO DEI DATI:
%     Schema del database, struttura JSON/XML, ontologie RDF, ecc.
%
%  LUNGHEZZA SUGGERITA: 20-35 pagine (tipicamente il capitolo piu' lungo).
% =============================================================================

% --- A. ARCHITETTURA ---
% Diagramma componenti + descrizione e giustificazione delle scelte architetturali.

% --- B. TECNOLOGIE ---
% Linguaggi, framework, librerie: versione, ruolo, motivazione della scelta.

% --- C. MODULI E FILE IMPORTANTI ---
% Struttura del codice. Usa \lstinputlisting o \begin{lstlisting}.

% --- D. ALGORITMI ---
% Algoritmi non banali: pseudocodice o codice con spiegazione passo-passo.
% Esempio:
% \begin{lstlisting}[language=Python, caption={Algoritmo XYZ}, label={lst:xyz}]
% def algoritmo(input):
%     ...
% \end{lstlisting}

% --- E. MODELLO DEI DATI ---
% Schema DB, strutture dati, ontologie, formati di scambio.

Scrivi qui il contenuto del capitolo 3 sull'architettura e l'implementazione.
